\chapter{Introduction}\label{introduction}

\section{Background and motivation}\label{background-and-motivation}

Online bipartite matching is the canonical model of irrevocable sequential allocation.
Resources are known in advance; requests arrive one at a time; each must be matched to an
available compatible resource --- or dropped --- immediately, before the rest of the input is
seen. It is the abstraction behind online advertising, ride-hailing dispatch, organ
exchange, and request routing in modern serving systems. Because decisions cannot be undone,
the difficulty lies entirely in committing well under uncertainty about the future.

A recent and influential idea is to give such an algorithm a \textbf{prediction} about the input
--- a hint about which resources will be contended, or how many requests of each kind will
arrive --- and to ask for two guarantees at once: \textbf{consistency}, meaning near-optimal
performance when the prediction is good, and \textbf{robustness}, meaning no worse than a
prediction-free algorithm when the prediction is wrong. For online matching specifically,
two families of such \emph{learning-augmented} algorithms have appeared: MinPredictedDegree,
which consumes a per-resource degree prediction, and a family of \emph{test-and-fallback}
schemes, which test a request-type prediction on a short prefix of arrivals before deciding
whether to trust it.

These algorithms come with attractive worst-case guarantees, yet a practitioner's
experience with them is oddly muted: on realistic, average-case workloads the sophisticated
predictions rarely seem to \emph{help}, while the algorithms' real value seems to be that they do
not \emph{hurt}. This gap between the worst-case promise and the average-case experience is the
puzzle that motivates this thesis. The algorithms had, moreover, only ever been studied \emph{in
isolation} --- each on its own input model, against its own notion of prediction error, and
largely through theory --- so there was no common ground on which to compare them, quantify
how much a prediction actually buys, or explain the gap. This thesis provides that common
ground and, ultimately, an explanation.

\section{Research objectives}\label{research-objectives}

The thesis is organized around three questions:

\begin{enumerate}
\def\labelenumi{\arabic{enumi}.}
\tightlist
\item
  \textbf{How do the learning-augmented online-matching algorithms actually compare}, head to
  head, under a single prediction-error model and on common inputs --- and how much does a
  prediction help, at what cost, on realistic data?
\item
  \textbf{What are the failure modes of the adaptive test-and-fallback mechanism} --- the cost of
  its test, the calibration of its accept/reject decision, and where it goes wrong?
\item
  \textbf{Why does the average-case experience differ so sharply from the worst-case promise} ---
  is the observed wall an artifact of particular algorithms and generators, or is it
  \emph{necessary}?
\end{enumerate}

The first two are answered experimentally; the third, theoretically.

\section{Contributions}\label{contributions}

\begin{itemize}
\item
  \textbf{A unified experimental benchmark} (Chapter 4) that places the advice-free baselines,
  MinPredictedDegree and its augmentations, and the test-and-fallback algorithms on one
  harness --- common graphs, a common structured prediction-error model, a common optimum, and
  confidence intervals --- the first head-to-head comparison of these families.
\item
  \textbf{A characterization of predictions as robustness insurance} (Chapters 4, 7): unguarded
  prediction-following crashes \emph{below} the advice-free baseline under bad predictions; two
  distinct mechanisms --- structural (the augmentations) and adaptive (test-and-fallback) ---
  restore safety with different consistency/robustness trade-offs; and the consistency
  upside is small on average-case inputs, so the value is downside protection. The picture
  is validated on synthetic graphs, six real-world graphs, and real request traces,
  including with a cheap, non-ML historical predictor.
\item
  \textbf{An empirical engagement with the order-error theory} (Chapter 5): MinPredictedDegree's
  loss is governed by a Kendall-\(\tau\) order error onto which several error models collapse.
  We credit Aamand--Chen--Indyk's Appendix D for the order-dependence itself and contribute a
  characterization of the known bound's looseness and saturation and of the right order
  measure --- not a rediscovery that order matters.
\item
  \textbf{The first empirical study of test-and-fallback} (Chapter 6): its testing cost, a
  threshold-calibration pathology in which a \emph{more accurate} test yields a \emph{worse} decision,
  its recalibration and the resolution limit that recalibration exposes, and a benchmark of
  the dynamic combiner revealing an irrevocability penalty that explains why matching
  requires \emph{test-then-commit}.
\item
  \textbf{An impossibility theorem} (Chapter 9): no test-and-fallback algorithm whose test
  inspects only a sublinear prefix can be both consistent and robust on strong-baseline
  instances, because the structure that makes the prefix distribution-test feasible is the
  structure that makes the baseline near-optimal. The proof reduces safe advice-use to
  \emph{tolerant} distribution testing and invokes its near-linear sample-complexity barrier;
  it holds for any decision rule, not merely the empirical threshold used in practice.
  \emph{(The theorem's full proof is complete in architecture --- its core lemma rigorous and its
  construction verified --- with one routine step and a final review outstanding; the thesis
  presents it with its status stated, as is appropriate for in-progress theoretical work.)}
\end{itemize}

Alongside these, the thesis documents (Chapter 8) the exploratory directions it pursued and
the negative results they returned --- an attempt to \emph{learn} the predictor for extra
performance, and an attempt to find a serving regime where predictions genuinely help ---
both of which reinforce the central finding and motivated the theorem.

\section{Thesis organization}\label{thesis-organization}

Chapter 2 surveys the background: online matching and its input models, the
learning-augmented paradigm, the specific prediction-based matching algorithms, and the
distribution-testing results the theory relies on. Chapter 3 fixes the model and
methodology and validates the experimental harness by reproducing a subset of the Borodin et
al.~study. Chapter 4 presents the unified benchmark and its four findings. Chapter 5
examines what governs the (small) loss --- order error --- and its relation to the known bound.
Chapter 6 studies the test-and-fallback mechanism in depth. Chapter 7 tests external
validity on real predictors and real graphs. Chapter 8 reports the exploratory directions
and negative results. Chapter 9 proves the impossibility theorem. Chapter 10 presents the
AI-inference serving case study. Chapter 11 concludes and discusses future work. The
recurring thesis, established experimentally and then proven necessary, is a single
sentence: \textbf{on average-case online matching, predictions are robustness insurance rather
than a performance lever --- and where you can test the advice, you do not need it.}
